\chapter{Introduction}\label{chap:intro}

\section{Motivation}\label{sec:motivation}
Network ossification is the process of a network growing
resistant to change.
This has become an issue at the transport layer, where
network middleboxes present a challenge to any IP
packets carrying a transport-layer protocol other than
the Transmission Control Protocol (TCP) or the User Datagram Protocol (UDP).
In addition to limiting innovation in
the form of new protocols, network ossification also limits
innovation within existing protocols such as TCP.
This limitation can be caused by implementations with bugs in rarely used features, or by network
middleboxes not supporting changes to the wire image~\cite{rfc8546} of a protocol.

Another barrier to new transport-layer protocols
is the coordination required to achieve widespread adoption.
Developers of server applications
will not spend time adding support for a new transport-layer protocol
if no applications are using it, and no applications
will use a new protocol if no servers support it.
If a fallback to TCP is always needed, then all applications
can avoid a great deal of complexity by only using TCP.

Together, network ossification and these deployment barriers
combine to make adoption of innovations at the transport layer difficult.

Transport Services (TAPS) is a series of IETF RFCs which aim
to tackle network ossification. TAPS presents
a common API which abstracts away protocol specific features, instead
allowing applications to specify which services they need from the
transport-layer protocol such as reliable data transfer or in-order delivery.
By creating a common interface, which is based on the application's
intent instead of concrete transport-layer protocols,
TAPS allows applications to be served by any
protocol that fulfills the requirements stated by the application.

Protocol fallback is built into TAPS,
requiring no extra complexity from applications. 
Applications written for QUIC can therefore remain compatible
with networks where only TCP is available, due to the
generic nature of the QUIC API.

\section{Problem Statement}\label{sec:prob-statement}
The TAPS working group concluded in 2025. Several open-source implementations of TAPS exist,
but none are actively maintained and none support QUIC. In addition, the TAPS standard provides no
guidance on how to map QUIC to the TAPS API.

This means that the only known implementation of TAPS, which supports both QUIC and
is actively maintained, is Apple's closed-source Network framework~\cite{networkFramework}.
This leaves two open questions: how to implement the current TAPS standard in practice,
and how QUIC features are expressed under the TAPS API.

\section{Research Questions}
The lack of an open, maintained TAPS implementation supporting QUIC,
combined with the absence of guidance on mapping QUIC to the TAPS API,
gives rise to two research questions:

\newcommand{\rqone}{What implicit architectural constraints does the TAPS specification impose on an implementation of a TAPS System?}
\newcommand{\rqtwo}{To what extent does the TAPS API expose the distinctive features of QUIC?}

\begin{enumerate}
    \item \rqone
    \item \rqtwo
\end{enumerate}

\section{Contributions}
The contributions of this thesis are CTaps, the insight that implementing
it has brought and a survey on how QUIC features can be mapped
to the TAPS API.
CTaps is an open-source C implementation of Transport Services with support
for QUIC, TCP and UDP. It is an exploratory work primarily
focused on how QUIC maps to the TAPS API and which
architectural concerns are encountered when implementing TAPS.

CTaps supports key QUIC features such as multistreaming, connection migration and session resumption.
Applications can utilize these features while remaining compatible with other
transport-layer protocols. CTaps also supports key TAPS features, such as candidate
gathering and candidate racing. This allows CTaps to race protocols and endpoints when
establishing connections.

\subsection{Scope}
CTaps covers the core TAPS API, with most functions fully implemented and parts
stubbed as no-ops. Rendezvous, Preconnection resolution and adding Endpoints to
established Connections are not supported.

Because CTaps is a research artifact, development has prioritized feature coverage over robustness
and error handling has not been as thoroughly tested as the core functionality.

\subsection{Source Code}
CTaps is open source and is available under the MIT license
on GitHub~\cite{cTaps}. Associated Doxygen documentation is available
as a static web page~\cite{ctapsDocs} hosted through GitHub pages.

\section{Thesis Structure}
\textbf{\Cref{chap:background}} covers the background necessary to understand
TAPS and CTaps.

\textbf{\Cref{chap:related-work}} surveys existing TAPS implementations,
comparing their features and development status to CTaps.

\textbf{\Cref{chap:design}} discusses the design of TAPS Systems
in general and CTaps in specific. Additionally, it discusses how well QUIC
features map through the TAPS API, weaknesses in the design of CTaps as well
as finally discussing security in a TAPS System.

\textbf{\Cref{chap:impl}} covers the internal design and implementation of
CTaps.

\textbf{\Cref{chap:eval}} compares CTaps with TCP, UDP and Picoquic clients, discussing
the advantages in terms of features and simplicity when utilizing CTaps.

\textbf{\Cref{chap:conc}} summarizes the conclusions of this thesis
and answers the posed research questions. It proposes future work
for CTaps and reflects on the broader significance of CTaps as a whole.
